\subsection{超边流量、串并行收益与划分尺度}
进一步说明为什么本问以共享张量超边而非普通边数参与初解。设一个输入张量 $x$ 大小 $b_x$，有三个消费算子 $u_1,u_2,u_3$。若三者处于同一 Task，题目数据路径至多需要该 Task 按核内策略读取该张量；若分在三个 Task，场景 A 的 Task 边界不保留私有缓存，至少有三次独立的消费机会。相对共同读取一次，静态重复输入代理由式\eqref{eq:q1-duplicate-input}给出 $2b_x$。但若用源到每个消费者的普通投影边表示，同一个张量可被计为三条边，简单地给每条边一个 $b_x$ 权重会得到 $3b_x$，和“重复读取增量”不是同一量；若再把跨子图读入 $3b_x$ 叠加到 $2b_x$，甚至重复统计。对内部张量又有生产侧写出和消费侧读入，场景 A 的真实开销取决于官方构图细节，不能机械套用输入张量的 $r_x-1$ 公式。因此超图权重只用于优先保留共享消费者，不作为正式 $D$ 的替代。

对于分量 $C$，可定义用于结构选择的三维轮廓
\begin{equation}\label{eq:q1-component-profile}
\mathbf q(C)=\bigl(W_M(C),W_V(C),R(C)\bigr),
\quad R(C)=\sum_{x:\,C(x)\cap C\ne\varnothing}b_x,
\end{equation}
其中 $R(C)$ 只表示潜在输入集合的字节规模。两分量没有原图依赖，分核不增加彼此的中间张量通信，但若它们共同消费原始输入，分到不同 Task 仍可能重复读入。这说明“弱连通分量完全独立”只能用于数据依赖和中间张量边的判断，不能宣称总输入搬运独立。轮廓中矩阵、向量工作量用于估计单核的两 Pipe 长尾；$R(C)$ 用于提醒共享输入代价。若只按 $W_M+W_V$ 做降序装箱，可能把一个 Cube 密集分量与另一 Cube 密集分量塞入同核，而让 Vector Pipe 和其它核心闲置。按式\eqref{eq:q1-pipe-load}比较添加后 $\max\{L_M,L_V\}$，可以把互补的 Cube/Vector 分量放在同核，利用两个计算 Pipe 的同时执行能力。

为了识别值得拆的主导分量，还要区分\emph{理论可并行余量}和\emph{实际可拆边界}。设完整分量 $C$ 的算子投影图在某拓扑位置切成 $C_1,C_2$，其理想负载下界从 $\max\{W_M(C),W_V(C)\}$ 可能降为 $\max\{\max_jW_M(C_j),\max_jW_V(C_j)\}$。令这两个数的差为 $G_{\rm pipe}$。切开后需要传递的超边集合记为 $\partial(C_1,C_2)$，其潜在 DDR 工作量至少与其中需要跨 Task 的不同张量有关。$G_{\rm pipe}$ 是\emph{释放并行的潜力}，不是预计 Makespan 的真实节省；边界流量、$100$/$1000$ 释放及局部溢出可抵消它。本文只在 $G_{\rm pipe}$ 足以提出合理候选、且商图保持无环时占用有限评价机会；如果整分量方案更快，保底机制选择整分量。由此，固定“每八算子一块”的规则被降为多种候选中的一类，而不作为普适划分尺度。

\subsection{边界移动的局部合法性与非局部代价}
设当前 $u\in V_a$，考虑把 $u$ 移到相邻子图 $V_b$。这一小动作需满足四层检查。第一，$V_a\setminus\{u\}$ 非空，且每个算子仍出现一次。第二，新的商图与核心顺序弧的并图 $H(p')$ 无环。第三，移动后的所有跨边数据都可按场景 A 插入合法的写出和读入事件，核内存储压力经官方排程处理后不违约。第四，候选必须在相同原图和同一核数下重新评价。前两项能用拓扑排序廉价地判定，后两项需要实际展开与仿真，不能仅凭原图局部邻接完成。尤其当 $u$ 是共享张量的生产者时，它的移动可能改变多个消费 Task 的 COPY 结构；看似局部的划分变更可对多核 DDR 在途集合产生非局部影响。

可用集合差而非边计数，预估移动后的输入超边变化。记 $R_x(p)=\{\pi(v):v\in C(x)\}$ 为张量 $x$ 的消费者 Task 集合，移动 $u$ 只可能改变包含 $u$ 的那些 $R_x$。代理差量为
\begin{equation}\label{eq:q1-local-hyperedge-delta}
\Delta\widehat D_{\rm input}(u:a\to b)=
\sum_{x:\,u\in C(x)}b_x\Bigl[
\bigl(|R_x(p')|-1\bigr)_+
-\bigl(|R_x(p)|-1\bigr)_+\Bigr].
\end{equation}
若 $a$ 中还有其它消费者而 $b$ 中已有消费者，则可能 $|R_x|$ 不变，移动算子并不会改变该输入的重复读取代理；若 $u$ 是 $a$ 中唯一消费者而 $b$ 中已有消费者，则 $|R_x|$ 减一。此增量可增量更新，因而适合从庞大候选集中筛出少量变体。但它仍不包含生产中间张量的写回、Pipe 排程和缓存溢出，所以不能代替式\eqref{eq:q1-added-bytes}。

\subsection{输出关键链与核心空闲的区别}
令 $\mathcal A(z)$ 为最终输出 $z$ 的算子祖先集合；若有多个输出，取最后完成的输出的祖先。某个算子或 Task 不在 $\mathcal A(z)$ 上，缩短它本身通常不能直接缩短当前实现时间，除非它在共享 Pipe 或 DDR 上抢占了关键祖先所需资源。这样可把候选分为两类：直接修改关键依赖路径的边界，或修改争用资源以间接帮助关键路径。单纯数核心忙碌百分比不能区分两类：一个核心全时忙碌但只做早早结束的非关键分支，与一个核心短时工作却在输出前形成最后释放点，对 Makespan 的贡献可能相反。因此我们同时观察最长 Task 链、跨核切换、Pipe 工作量和共享 DDR 时窗。

若某图有一个核心长期空闲，也不能简单强制把任务迁到它。考虑迁移一个 Task $r$ 的估计净效应：释放原核心的计算负载 $\Delta C_{\rm old}$，增加新核心的负载 $\Delta C_{\rm new}$，并改变进入与离开 $r$ 的 Task 边跨核状态。即使 $\Delta C_{\rm new}=0$、原核上有明显尾部队列，新增的关键链跨核释放仍可能使输出晚于原方案。一个必要的结构检查是：迁移后所有受影响的前驱和后继能否在合法顺序中释放；一个性能检查是完整仿真得到的 $T(p')<T(p)$。计算空闲只是候选触发信号，而不是接受条件。对问题一尤其如此，因为同核 Task 边界本就经 DDR，不存在“迁到同核就免传输”的捷径；迁核主要改变释放等待和并行排程。

\subsection{多核加速比的两个上限直觉及其局限}
若把全部计算视作一条完全可分的单类 Pipe，理想速度增益似乎至多为核数 $k$；但题目以\emph{启发式整图单核}结果作分母，该基准可能因整图规模与缓存压力而不是理论最优，故实测比值超过 $k$ 不构成逻辑错误。另一方面，若只从固定计算链看，$L_{\rm op}$ 给出无论核数多少都无法突破的时间下界，因此对某个方案的参考单核时间 $T^0_{i,1}$ 有形式上的上限 $S_{i,k}\le T^0_{i,1}/L_{\rm op}$。该上限往往很松，因为没有算入同步、DDR 和容量。本文不拿它预测成绩，只用来说明当关键串行链占主导时增加核心不可能无限加速，以及为何需要在分量和关键链之间平衡切图。

更实际的平均指标也不能由“每个图最优地利用五核”推断。若图 $i$ 的 $T^0_{i,1}$ 很小、受固定 Task 释放主导，则 $S_{i,5}$ 可能低于其它图；若图 $j$ 的整图核内溢出异常大，切成多核后即使新增通信仍可能得到超过五倍的比值。式\eqref{eq:mean-speedup-derivation}给每张图相同权重，要求算法在不同结构之间稳定，而不是只针对总周期占比最大的几张大图取得很好的成绩。我们在报告中并列列出大图周期贡献与百图逐图均值，目的不是制造两个“官方分数”，而是说明总目标与统计呈现的不同关注点。
